
\subsubsection{4.1 Proof-of-Concept Scope}

Validation uses 100 models (50 ResNet-50, 50 ViT-Base) with synthetic accuracy labels sampled from realistic ImageNet top-1 distributions: mean 76.1\%, std 2.3\% for ResNet-50; mean 81.8\%, std 1.7\% for ViT-Base. Synthetic labels isolate mechanism validation---verifying that operation encoders produce distinct representations, contrastive projection aligns tasks while preserving structure, and GNN adds topology signal---from empirical property prediction. Full-scale validation with 400 models across 4 architectures and real accuracy labels is planned for future work.

\subsubsection{4.2 Model Zoo}

Models were collected from HuggingFace Hub: 50 ResNet-50 variants and 50 ViT-Base variants, all nominally trained on ImageNet-1K classification. Models span diverse training configurations including different random seeds, data augmentation strategies, and optimization schedules.

\subsubsection{4.3 Evaluation Protocol}

Train/validation/test split: 70/15/15 by count, stratified by architecture. Training set: 70 models (35 ResNet-50, 35 ViT-Base). Validation set: 15 models. Test set: 15 models.

Primary metric: Spearman rank correlation $\rho$ between predicted and actual (synthetic) ImageNet top-1 accuracy for $ResNet\rightarrow ViT$ cross-architecture transfer. Gate threshold: $\rho \geq$ 0.65, representing 35\% gap closure between SNE cross-architecture ($\rho$ = 0.54) and within-architecture performance ($\rho$ = 0.81).

Statistical validation: Permutation test with 1000 iterations. For each iteration, architecture labels are randomly shuffled and SNE baseline is retrained. CAPE's $\Delta\rho = \rho_{CAPE} - \rho_{SNE}$ is compared against the null distribution. If p < 0.05, improvement is statistically significant.

\subsubsection{4.4 Ablation Study}

Four variants isolate component contributions:

\begin{enumerate}
\item \textbf{SNE Baseline:} Set-encoding with no operation-specific encoders, no contrastive projection, no GNN.
\item \textbf{Operation-Only:} Add modular operation encoders but no contrastive projection or GNN.
\item \textbf{Op+Contrastive:} Add contrastive projection to operation encoders but no GNN residual.
\item \textbf{Full CAPE:} Add architecture GNN with learnable residual weight $\alpha$.
\end{enumerate}

All variants use identical hyperparameters and training data. Each trains independently.

\subsubsection{4.5 Diagnostic Metrics}

Three pre-registered falsifiers verify mechanisms:

\begin{enumerate}
\item \textbf{Operation Encoder Distinctiveness:} Cosine similarity between z\_conv and z\_attn. Threshold: < 0.95. If exceeded, modular encoding fails.
\end{enumerate}

\begin{enumerate}
\item \textbf{Contrastive Alignment Quality:} Intra-architecture variance in z\_proj space. Threshold: $\geq$ 0.1. If below, contrastive alignment collapses structure.
\end{enumerate}

\begin{enumerate}
\item \textbf{GNN Residual Contribution:} Learned residual weight $\alpha$. Threshold: > 0.1. If below or if adding z\_arch decreases performance, GNN is ineffective.
\end{enumerate}

\subsubsection{4.6 Computational Environment}

Training executed on $5\times$ NVIDIA H100 NVL GPUs (95GB memory each). PyTorch 2.7.1 with CUDA 11.8. PyTorch Geometric 2.8.0 for GNN operations. Total training time: 45 minutes for 10 epochs at proof-of-concept scale.
